% Fragmento público generado desde propietarios íntegros.
% Residencia: 52.5.1.
% Fuente: ../incoming/radion_monografia_20260827/manuscrito/sections/06_interior_helice.tex:1-108
\noindent Projective geometry distinguishes the finite action from hyperbolic depth and return memory.\par\medskip

\subsubsection{Hyperbolic Interior of Hilbert--Beltrami--Klein}

\paragraph{Projective normalization on the Klein disk}

The affine application
\[
u=\frac{Q}{a_S},\qquad v=\frac{P}{b_S}
\]
carries the action ellipse to the unit disk
\[
D=\{(u,v):u^2+v^2<1\}.
\]
On each chord, the Hilbert distance is defined through the cross-ratio
of its boundary endpoints. The resulting infinitesimal metric is the
Beltrami--Klein form
\begin{equation}
\diff s_K^2=
\frac{(1-r^2)(\diff u^2+\diff v^2)+(u\diff u+v\diff v)^2}
{(1-r^2)^2},
\qquad r^2=u^2+v^2.
\label{eq:klein-metric}
\end{equation}
With this normalization, the Gaussian curvature is constant:
\begin{equation}
K=-1.
\label{eq:klein-curvature}
\end{equation}
The geodesics are Euclidean chords, but their hyperbolic length diverges as
they approach the boundary.

\paragraph{Finiteness of the boundary and infinitude of the hyperbolic depth}

The metric determinant of \eqref{eq:klein-metric} produce
\begin{equation}
\diff A_K=\frac{\diff u\,\diff v}{(1-r^2)^{3/2}}.
\label{eq:klein-area-density}
\end{equation}
The area of the subdisk \(r<R<1\) is
\begin{align}
A_K(R)
&=\int_0^{2\pi}\int_0^R
\frac{r\,\diff r\,\diff\theta}{(1-r^2)^{3/2}}\\
&=2\pi\left(\frac1{\sqrt{1-R^2}}-1\right).
\label{eq:klein-area-R}
\end{align}
\Needspace{4\baselineskip}
Therefore,
\[
\lim_{R\to1^-}A_K(R)=+\infty.
\]
Simultaneously, the same boundary encloses the symplectic area.
\[
\pi a_Sb_S=2\pi\hret=\hfull<\infty.
\]

\begin{theorem}[Finiteness of Action and Infinitude of Depth]
The radion ellipse has a finite turn action \(\hfull\), whereas
its interior, equipped with the Hilbert--Beltrami--Klein metric, has infinite
hyperbolic area:
\begin{equation}
\boxed{
\operatorname{Area}_{\mathrm{simp}}(\partial\mathbb E_S)=\hfull<\infty,
\qquad
\operatorname{Area}_{K}(\mathbb E_S)=+\infty.}
\label{eq:finite-infinite}
\end{equation}
\end{theorem}

There is no contradiction: \(\operatorname{Area}_{\mathrm{simp}}\) and
\(\operatorname{Area}_K\) are different measures. The first counts oriented action
in phase; the second measures interior projective depth. The thesis
\enquote{finite boundary, infinite interior} is now an equality and a limit,
not a metaphor.

\paragraph{Focal Coordinates Inside Klein}

In the normalized disk, the foci are located at \((\pm e_f,0)\). The distance from the center to a focus is
\begin{equation}
u_f=\operatorname{artanh}e_f
=\operatorname{arcosh}(e^\eta)
=0.813382331175342333760889731088228\ldots.
\label{eq:klein-focal}
\end{equation}
The focal--focal distance is
\[
d_K(F_-,F_+)=2u_f
=1.62676466235068466752177946217646\ldots.
\]
We have
\begin{equation}
e_f=\tanh u_f,\qquad
e^{-\eta}=\operatorname{sech}u_f,\qquad
\eta=\log\cosh u_f.
\label{eq:focal-hyperbolic-chain}
\end{equation}
The region up to the focal radius has area
\[
A_K(e_f)=2\pi(e^\eta-1)
=2.19559620884061869541206115980360\ldots.
\]

All open ellipses are projectively isometric as Hilbert domains. Curvature
\(-1\) alone therefore does not retain eccentricity. The HMT form
resides in the joint structure: the distinguished affine chart, symplectic
form, sheet labeling, and boundary parametrization.
